\par\Needspace{6\baselineskip}\medskip
\begin{center}\hypertarget{native-235-article170}{}\pdfbookmark[2]{170.}{article-170}170.\end{center}
In der ersten Hypothese müssen wir die unkorrigierten Zeiten $t,t',t'',t'''$ verwenden, weil sich die Abstände von der Erde noch nicht berechnen lassen. Sobald aber Näherungswerte der Größen $x',x''$ bekannt sind, können wir auch diese Abstände annähernd bestimmen. Da die Formeln für $\rho$ und $\rho'''$ hier etwas komplizierter sind, empfiehlt es sich, die Berechnung der Zeitkorrektur aufzuschieben, bis die Abstände hinreichend genau bestimmt sind. So erübrigt sich eine Wiederholung der \SourceLeaf{236}Rechnung. Daher legen wir dieser Rechnung zweckmäßigerweise die Werte von $x',x''$ zugrunde, die sich aus der vorletzten Hypothese ergeben haben. Erst die letzte Hypothese geht dann von den korrigierten Zeiten und Werten der Größen $P',P'',Q',Q''$ aus. Hierfür dienen die folgenden Formeln.
\[
\tag*{VII.}
\rho'=x'-R'\cos\delta'
\]
\[
\tag*{VIII.}
\rho''=x''-R''\cos\delta''
\]
\[
\tag*{IX.}
\begin{aligned}
\rho\cos\beta={}&-R\cos B\cos(a-l)\\
&+\frac{1+P'}{P'\left(1+\dfrac{Q'}{r'^3}\right)}\bigl(\rho'\cos\beta'\cos(a'-a)+R'\cos B'\cos(l'-a)\bigr)\\
&-\frac{1}{P'}\bigl(\rho''\cos\beta''\cos(a''-a)+R''\cos B''\cos(l''-a)\bigr)
\end{aligned}
\]
\[
\tag*{X.}
\begin{aligned}
\rho\sin\beta={}&-R\sin B\\
&+\frac{1+P'}{P'\left(1+\dfrac{Q'}{r'^3}\right)}\bigl(\rho'\sin\beta'+R'\sin B'\bigr)\\
&-\frac{1}{P'}\bigl(\rho''\sin\beta''+R''\sin B''\bigr)
\end{aligned}
\]
\[
\tag*{XI.}
\begin{aligned}
\rho'''\cos\beta'''={}&-R'''\cos B'''\cos(a'''-l''')\\
&+\frac{1+P''}{P''\left(1+\dfrac{Q''}{r''^3}\right)}\bigl(\rho''\cos\beta''\cos(a'''-a'')\\
&\qquad{}+R''\cos B''\cos(a'''-l'')\bigr)\\
&-\frac{1}{P''}\bigl(\rho'\cos\beta'\cos(a'''-a')+R'\cos B'\cos(a'''-l')\bigr)
\end{aligned}
\]
\[
\tag*{XII.}
\begin{aligned}
\rho'''\sin\beta'''={}&-R'''\sin B'''\\
&+\frac{1+P''}{P''\left(1+\dfrac{Q''}{r''^3}\right)}\bigl(\rho''\sin\beta''+R''\sin B''\bigr)\\
&-\frac{1}{P''}\bigl(\rho'\sin\beta'+R'\sin B'\bigr).
\end{aligned}
\]
Die Formeln IX—XII lassen sich ohne Mühe aus den Gleichungen 1, 2, 3 in Artikel 112 ableiten, wenn man die dort verwendeten Zeichen in die hier verwendeten überträgt. Die Formeln werden offenbar viel einfacher, wenn $B,B',B''$ verschwinden. Aus der Verbindung der Formeln IX und X ergeben sich sowohl $\rho$ als auch $\beta$ und ebenso aus XI und XII neben $\rho'''$ auch $\beta'''$. Sind die beobachteten Breiten (die nicht in die Rechnung eingehen) gegeben, so zeigt der Vergleich mit ihnen, mit welcher Genauigkeit sich die Breiten der ersten und letzten Beobachtung durch die sechs an die übrigen Angaben angepassten Elemente darstellen lassen.
